\documentclass[11pt,a4paper]{article}
\usepackage[margin=2.5cm]{geometry}
\usepackage{enumitem}
\usepackage{booktabs}
\usepackage{hyperref}
\usepackage{titlesec}
\usepackage{fancyhdr}
\usepackage{xcolor}
\usepackage{tabularx}

\hypersetup{
    colorlinks=true,
    linkcolor=blue!60!black,
    urlcolor=blue!60!black
}

\titleformat{\section}{\large\bfseries}{}{0em}{}
\titleformat{\subsection}{\normalsize\bfseries}{}{0em}{}

\pagestyle{fancy}
\fancyhf{}
\rhead{\small\thepage}
\lhead{\small MA International Relations --- Thesis Guidelines}
\renewcommand{\headrulewidth}{0.4pt}

\setlength{\parindent}{0pt}
\setlength{\parskip}{0.8em}

\begin{document}

\begin{center}
{\Large\bfseries MAIR THESIS GUIDELINES}\\[0.3em]
{\normalsize Leiden University $\cdot$ Faculty of Humanities}\\[0.3em]
{\small Based on MAIR Thesis Supervision Guidelines, last updated December 05, 2025}\\[1.5em]
\end{center}

\begin{quote}
\small\textit{This document is adapted from the MAIR Thesis Supervision Guidelines for student reference. Some content intended for supervisors and second readers has been removed. For the most current information, always consult your supervisor and check Brightspace.}
\end{quote}

\bigskip
\hrule
\bigskip

\section{1. Supervisor Assignment}

In the second block of your first semester, you submit a literature review on the basis of which you are assigned to a supervisor by your MAIR specialization---i.e., Culture and Politics, European Union Studies, Global Conflict in the Modern Era, Global Order in Historical Perspective, and Global Political Economy. You may indicate a preference for supervisors whose research interests match your own.

Note that changes of supervisor are generally not facilitated except in extraordinary circumstances. A change of topic is generally not sufficient cause for a change of supervisor.

\section{2. Supervision}

Supervision is expected to start upon assignment. Your supervisor will meet with you shortly after assignment to help you set goals, including for the upcoming winter or summer break (during which supervisors are generally not available for supervision).

\textbf{Minimum meetings:} You should meet with your supervisor at least four times:

\begin{enumerate}
    \item \textbf{Introductory meeting:} Discussion of the choice and scope of the thesis topic, the research question, the literature, the source materials, the general approach, and the timetable for submissions.
    \item \textbf{Introduction/literature review/research design:} Discussion following submission of these sections.
    \item \textbf{Empirical chapter:} Discussion of an empirical chapter during the intermediate stage.
    \item \textbf{Final evaluation:} Discussion of the final evaluation.
\end{enumerate}

\textbf{Feedback:} Your supervisor will provide oral or written feedback on the introduction/literature review/research design and at least one subsequent chapter. Supervisors are allowed but not required to read and provide feedback on an entire draft. Your supervisor will communicate their supervision style during the first meeting.

\textbf{Writing support:} If you are struggling with writing in English, contact the faculty's \href{https://www.student.universiteitleiden.nl/en/vr/humanities/writing-lab}{Writing Lab}.

\textbf{Important:} Because there is no supervision during the winter and summer breaks, theses cannot be written in their entirety over the break. The student should actively consult their supervisor during the research process---a thesis cannot be submitted for evaluation that has not been overseen by a supervisor.

\section{3. Research Ethics}

You and your supervisor should discuss any potential ethical implications of your research during the initial meeting or when they first arise. Ethical implications should always be considered when the research involves human participants or the use of individually identifiable data. Particular care should be taken with vulnerable groups, such as asylum seekers or citizens of authoritarian regimes. \textbf{The program will not allow any thesis research that has the potential to endanger the researcher or participants.}

\subsection{When conducting interviews:}

\begin{itemize}
    \item Acquire \textbf{informed, voluntary consent} from all participants by providing an information sheet that explains the nature, aims, and implications of the research project, the nature of participation, and any other considerations that might reasonably influence willingness to participate.
    \item Protect the \textbf{privacy and confidentiality} of participants throughout the research process. Information that may identify individual persons should not be used unless the person has expressly agreed.
    \item Use a \textbf{consent form} that includes acknowledgement that participants have read and understood the information sheet, and details whether they consent to being identified or prefer anonymity.
    \item Respect the right of individuals to \textbf{refuse to participate} or withdraw their consent at any stage without prejudice.
\end{itemize}

For additional information, see the Code of Ethics available on the \href{https://scdenney.github.io/thesis-supervision/ethics/}{Ethics \& Policies page}.

\section{4. Thesis Deadlines and Submission}

There is a single deadline to submit the thesis each semester. If the deadline is not met, the next opportunity is the following semester.

\begin{center}
\begin{tabularx}{\textwidth}{ll}
\toprule
\textbf{Semester} & \textbf{Deadline} \\
\midrule
Spring 2026 & Friday, June 5, 2026 \\
Fall 2026 & TBD (typically first Friday of December) \\
\bottomrule
\end{tabularx}
\end{center}

Submit your thesis by email to your supervisor with CC to the second reader. Ask for confirmation that your thesis is received and submitted. If you do not receive confirmation, follow up after one week.

\subsection{Extensions}

If you need to extend the deadline, you must discuss this with your supervisor and second reader---\textbf{both} need to approve the extension. Please consider that:

\begin{itemize}
    \item Extension requests must be discussed well before the deadline
    \item Extensions can result in a later graduation date
    \item Extensions may prevent participation in the graduation ceremony
\end{itemize}

If the extension request is not approved, contact \href{mailto:stucomair@hum.leidenuniv.nl}{stucomair@hum.leidenuniv.nl} and consult the Board of Examiners, who base their decision on a hindrance statement from the student counsellors.

\subsection{Resubmission}

Students who fail their thesis are allowed a single retake. Students who resubmit their revised thesis by the last workday of January (for the December deadline) or August (for the June deadline) can still graduate within the same month.

\section{5. Style and Formatting}

\begin{itemize}
    \item \textbf{Maximum word count:} 15,000 words, including all elements (notes, bibliography, and appendices). This is a \textbf{hard maximum} and does not allow for a 10\% margin.
    \item One reference style must be used consistently.
    \item The thesis should:
    \begin{itemize}
        \item be written in English;
        \item be written in Word;
        \item use 1.5 line spacing, standard margins, and a standard size 12 font;
        \item include the student's name, email, student number, and word count on the title page;
        \item include the student number on all subsequent pages;
        \item be proofread for spelling and language errors.
    \end{itemize}
\end{itemize}

\section{6. Assessment}

Both the first and second readers will assess the thesis based on the criteria set out below.

Please note that the research \textit{process} is a necessary component for evaluation. While the thesis is an independent project, it cannot be submitted for evaluation without having been overseen by a supervisor.

\subsection{Quality Criteria}

\begin{itemize}
    \item The thesis should include a clearly formulated research question.
    \item The thesis should include a critical and analytical report on existing academic debates based on secondary literature.
    \item The thesis should include an original contribution to academic debates. In principle, the thesis should be of sufficient quality (possibly following some modifications) to be published in an academic journal.
    \item The thesis should be based on primary sources, if applicable and appropriate.
    \item The thesis should display the application of concepts and research methods.
    \item The thesis should display clear structure and proper use of language.
    \item \textbf{Fraud and plagiarism are `knockout criteria.'} Suspicions thereof must be reported to the Exam Committee.
\end{itemize}

\subsection{Fit with Specialization}

The thesis should fit your specialization. The thesis is envisaged as the culmination of studying a particular set of core courses and electives. You should write a thesis that builds upon, and makes use of, the knowledge, topics, methods, and/or skills gained from following your specific specialization.

\subsection{Assessment Criteria}

\textbf{Knowledge and Insight:} The research question is based on a problem that reflects insight into the key discussions and methods of the field; clarity, relevance, and definition of the problem; embedding in the existing literature; originality.

\textbf{Application of Knowledge and Insight:} Critical analysis of primary sources; use of complex concepts and effective research methods; description and justification of the adopted method; application of knowledge in broader or multidisciplinary contexts.

\textbf{Reaching Conclusions:} Logical and consistent reasoning; well-founded conclusions; degree to which the thesis question is answered; connection to other and future research; social and ethical aspects; critical reflection on own role as researcher.

\textbf{Communication:} Language use (readability, style, grammar, terminology); structure and layout; correct use of citations and bibliography.

\textbf{Learning Skills (Process):} Degree of independence; planning and time management; handling feedback from supervisors; if applicable, participation in thesis group.

\section{7. Grade Descriptors}

\begin{center}
\small
\begin{tabularx}{\textwidth}{llX}
\toprule
\textbf{Grade} & \textbf{Class} & \textbf{Description} \\
\midrule
9--10 & A+ Distinction & Outstanding; excellent understanding of issues and methodologies; original, independent thinking; rigorous argument with wide range of sources. At 10, could not be bettered at the MA level. \\
\midrule
8--8.9 & A Merit & Excellent understanding; independent thought; strong, well-organized argument using a wide range of sources. \\
\midrule
7--7.9 & B Merit & Good to very good; most but not necessarily all of the above. \\
\midrule
6--6.9 & C Pass & Satisfactory understanding; reasonable argument with standard range of sources. Some shortcomings but no fundamental errors. \\
\midrule
5.1--5.9 & & \textit{The faculty does not issue grades in this range.} \\
\midrule
3--5.0 & F Fail & Inadequate understanding; substantial omissions or irrelevant material. Poorly conceived. \\
\midrule
2--2.9 & F Fail & An attempt to answer, but without significant grasp of material or appropriate skills. \\
\midrule
0--1.9 & U & No answer; totally irrelevant, fundamentally wrong, or plagiarized. \\
\bottomrule
\end{tabularx}
\end{center}

\section{8. Plagiarism}

The unacknowledged reproduction of the work, words, or ideas of another person constitutes plagiarism. Students should cite their sources and paraphrase---never `copy' or moderately edit original text. Any work found to be plagiarised or to contain plagiarism will receive, at the very least, a mark of zero. The maximum penalty is expulsion from the programme.

See the \href{https://www.organisatiegids.universiteitleiden.nl/en/regulations/general/plagiarism}{Regulations on Plagiarism} and the \href{https://www.organisatiegids.universiteitleiden.nl/en/regulations/humanities/guidelines-for-the-use-of-genai-in-assessment}{Guidelines for the Use of GenAI in Assessment} established by the Faculty of Humanities.

\end{document}
